\documentclass[12pt]{article}
\usepackage[margin=0.7in]{geometry}

% \usepackage[margin=0.8in]{geometry}

% \usepackage{newpxtext}
% \usepackage{newpxmath}

\usepackage[tt=false]{libertine}

% \usepackage{newtxtext}
% \usepackage{newtxmath}

\usepackage[stretch=10]{microtype}

\linespread{1.05}

\usepackage{enumitem}
\setlist{nosep}

\usepackage[hidelinks, linktoc=all]{hyperref}

% Answers to the review questions. In the PDF this typesets the answer as an
% indented block right after the question; on the web the same environment is
% turned into a click-to-reveal dropdown (see web-pipeline/qa-details.lua).
% Pandoc does not implement \NewDocumentEnvironment, so it leaves the environment
% as a Div the filter can restructure, while LaTeX honours the definition here.
\NewDocumentEnvironment{answer}{}
  {\par\smallskip\noindent\ignorespaces}
  {\par\smallskip}

\title{Ethics}
\date{}
\author{}


\begin{document}

\section{Deontology}
Deontology is a duty-based ethical theory asserting that the morality of an action is based on whether the action itself is right or wrong under a series of rules, rather than based on the consequences.

\subsection{Categorical Imperative}
According to Kant, the Categorical Imperative (CI) is the central, unconditional moral law that dictates actions must be based on duty rather than desire.

\noindent\textbf{Hypothetical Imperative (Desire-based)}
"\textit{If you want X, do Y}". These rules are based entirely on our personal inclinations or desires.

\noindent In contrast, CIs are \textbf{reason-based}.
"\textit{I should do X whether or not I want anything else}".
These are moral commands based on reason alone, completely ignoring what you personally want.
In other words, the rule is followed purely out of a sense of duty, rather than a desire for a reward, this action holds genuine moral goodness.

\subsection{Principle of Universalisability}
This is the first formulation of Kant's Categorical Imperative.

\noindent PoU says "you should act only on principles that you would be willing to have absolutely everyone else act on as well".
In other words, an action is morally right if and only if it is universalisable.

\noindent\textbf{The Universalisability Test}
\begin{enumerate}
    \item Imagine a hypothetical world in which absolutely everyone performs the action you are about to take.
    \item Confirm that in this hypothetical world, both of the following statements are true:
    \begin{itemize}
        \item The goal of the action can still actually be achieved in such a world.
        \item Nothing essential to the agent’s (your) will would be endangered if everyone in the world performed this action.
    \end{itemize}
\end{enumerate}
If both of those statements are true, the action is universalizable, which means it is morally right. 
If either statement is false, the action is a contradiction and is morally wrong.

\noindent The fundamental thought behind this principle is that acting morally strictly requires being impartial, and not making an exception for themselves.

\noindent\textbf{Example - Lying}

\noindent Imagine there is something very expensive you want to buy, but you lack the money.
You have a wealthy acquaintance who you know will give you a loan if you simply ask.
This loan would fully cover the cost of the item you want.
You know for a fact that if you never repaid the loan, you would not suffer any punishment or retaliation from this acquaintance or anyone else.

\noindent Applying the test, this action is immoral (fails statement 1).
If absolutely everyone in the world constantly made false promises to get loans, the entire social concept of a "promise" or a "loan" would completely collapse. No one would ever believe anyone else, and consequently, no one would ever lend out money.
Because your goal (getting the loan by making a promise) is logically impossible to achieve in a world where everyone acts the same way, the action contradicts itself.

\noindent\textbf{Example - Not helping}

\noindent Imagine this maxim: "I will never help anyone, and I expect no one to ever help me."
Applying the test, this world can technically exist (passes statement 1), although it would be a cold, cruel world.

\noindent However, Kant points out a biological and rational fact about human beings: we are not self-sufficient.
As a rational human being, you inevitably will that somebody helps you when you are in danger.
This contradicts the will that nobody helps you (fails statement 2).
Therefore, we have a duty to help others.

\vspace{1em}
\noindent\textbf{Problem with Universalisability Test}
The Universalizability test is highly vulnerable to how you phrase your maxim, leading to "False Positives.
If your maxim is "I will rob a bank to get money," it fails the test (if everyone robs banks, the concept of money collapses).
But if you phrase it very specifically: "I will rob a bank only on Tuesdays in July while wearing a silly hat", it might actually pass the test.
If everyone in the world who wore a silly hat on a Tuesday in July robbed a bank, the banking system would likely survive.
The test can be tricked by overly specific wording.

\subsection{Principle of Humanity}
This is the second formulation of Kant's Categorical Imperative.
An action is morally right if and only if it treats human beings as ends in themselves, and never merely as means to an end.
\begin{enumerate}
    \item Never use others as \textbf{mere} means to an end.
    To use someone as a "mere means" is to use them strictly as a tool for purposes that are not their own purposes.
    For example, if you make a lying promise to get a loan, you are using that person purely as a financial tool.
    However, Kant notes that it is morally permissible to use others for purposes that they actively share with you  (e.g., you use a barista to get coffee, but they share the purpose of completing that transaction).
    \item Always treat others as ends in themselves.
    To treat someone as an "end in themselves" means treating them with respect and recognizing them as a rational, autonomous being with intrinsic value and dignity.
    For example, telling the truth allows others to use their own rationality to decide if they want to help you.
\end{enumerate}
"Human beings" refer to all rational, autonomous agents.
\begin{itemize}
    \item Rationality: The power to reason and to use that reason to govern your own life.
    \item Autonomy: The capacity to govern your own life by making your own free choices.
\end{itemize}
All rational, autonomous beings are intrinsically valuable, and they all possess equal intrinsic value.

\noindent\textbf{Implications on our moral obligations}
\begin{itemize}
    \item Negative Obligations (Constraints):
    We are morally required to refrain from diminishing other people's rationality and autonomy, as well as our own. 
    This creates strict obligations not to lie, steal, enslave, rape, or murder, because all of these actions involve using people for your own purposes.
    \item Positive Obligations:
    We are also required to actively promote people's rationality and autonomy. 
    For example, you have a moral obligation to perform easy rescues (like saving a drowning toddler) because death destroys a rational autonomous being. 
    It also means we have an obligation to help reduce hunger, great poverty, and powerlessness.
\end{itemize}

\subsection{Perfect vs Imperfect Duties}
Kant divided our moral obligations (generated by the CI)
into Perfect and Imperfect duties.

\noindent\textbf{Perfect duties} are mandatory, strict rules that must always be followed with no exceptions.
Violating them creates a "contradiction in conception" (Statement 1, logically impossible to universalize).
They are usually negative duties:
\begin{itemize}
    \item Do not lie
    \item Do not steal
    \item Do not murder
    \item Do not break promises
\end{itemize}

\noindent\textbf{Imperfect duties} are meritorious, general moral obligations we should adopt as lifelong goals,
but you have the flexibility to choose when, how, and toward whom you fulfill them.
Violating them is not a logical contradiction, but a "contradiction in will" (Statement 2, we would not want to live in a world where everyone acted that way).
They are usually positive duties:
\begin{itemize}
    \item To help others in need
    \item To develop one's own talents
\end{itemize}

\noindent\textbf{Conflict} If a perfect duty conflicts with an imperfect duty,
we should prioritise the perfect duty because it is a strict obligation.
For example, one must not steal (perfect) even to help someone in need (imperfect).
The ends never justify the means.

Because Perfect Duties are negative constraints that protect basic human autonomy (Contradiction in Conception), while
Imperfect Duties are positive goals (Contradiction in Will).
You cannot actively violate someone's basic autonomy (break a perfect duty) just to achieve a positive goal (fulfill an imperfect duty).

\vspace{1em}
\noindent\textbf{Perfect/Imperfect vs Obligatory/Supererogatory}
Imperfect duties (like helping others or developing talents) are not to be confused with supererogatory actions (like donating large sum to charity).
An imperfect duty is still a strict moral requirement. You must adopt the goal of helping people. The only "imperfect" part is that you have the flexibility to choose when, where, and how you fulfill it.

\subsection{2 Objections Against Deontology}
\subsubsection{Inquiring Murderer}
\textbf{Scenario} Imagine someone you love is fleeing from a murderer and tells you where she is hiding. The murderer comes to your door and inquires about her location. Should you lie?

\noindent Assumptions:
\begin{itemize}
    \item The inquiring murderer is no threat to you personally; he will not attack you if he finds out you lied to him.
    \item If you don’t lie, the murderer will most probably find and kill the intended victim.
    \item If you do lie, the murderer will most probably not find the intended victim.
\end{itemize}

\noindent According to Kant's CI, we should never lie, as lying violates the CI.

\noindent Benjamin Constant's claim is that it would not be immoral to lie to the inquiring murderer. In fact, lying here may be morally required
to save a life.

\noindent\textbf{Kantian defence} Although the CI is a moral absolute,
it may be more flexible than even Kant himself realized.
The CI can demonstrate that it would be morally right to lie in some situations.
Specifically, they argue that lying to the murderer can actually pass the Universalizability test.
\begin{itemize}
    \item The goal of this action can be achieved in this world. 
    In such a world, an inquiring murderer wouldn’t know for sure whether anyone was lying to him. 
    For all he knows, they really have no idea. 
    Contrast this with a world in which everyone makes a lying promise to repay a debt; in that world, as soon as someone promises to repay a debt, you know for sure they are lying.
    \item Nothing essential to the agent’s will would be endangered. 
    In fact, it is protected because more innocent lives would be saved. 
    Your will could enjoy the protection of others who would lie to save you from being killed.
\end{itemize}


\subsubsection{Animal Rights}
Kant's Principle of Humanity specifically states that we must treat beings with rationality and autonomy as ends-in-themselves.
However, this excludes beings lacking rationality, like infants, individuals with severe mental disabilities, and non-human animals from receiving direct moral consideration.

\noindent This is in contrast from utilitarian approach that asks "\textit{Can they suffer?}".

\noindent For example, the statement "\textit{animal cruelty is perfectly fine}" goes against basic human intuition.
Kant created a workaround: He argued that we do not have direct duties to animals, but rather indirect duties to them because of how our actions affect us.
\begin{enumerate}
    \item \textbf{Premise 1}: Cruelty to animals makes you more likely to be cruel toward rational, autonomous beings (other humans). Kant believed that abusing animals degrades our natural sense of empathy and makes us more cruel overall.
    \item \textbf{Premise 2}: It is immoral to act in a way that makes you more likely to be cruel toward rational, autonomous beings.
    \item \textbf{Conclusion}: Therefore, cruelty to animals is immoral.
\end{enumerate}

\noindent Ultimately, our duties toward non-rational animals and nature are really just duties to respect our own humanity, which requires us to cultivate virtues like kindness and gratitude.

\subsection{Agent-Centered vs Victim-Centered Deontology}
\textbf{Agent-centered Deontology} focuses on the moral agent’s actions and intentions.
Agents have special, "agent-relative" obligations to avoid doing wrong, even if doing so results in more harm overall.
For example, parents have agent-relative obligations to their own children. Furthermore, parents are granted the agent-relative permission to save their own children, even if it comes at the expense of other children not being saved.

\noindent Doctrine of double effect: an action with both a good effect and a foreseen, unintended bad side effect may be morally permissible if the intention is solely to achieve the good result.
For example, in the Trolley Problem, if your explicit intention is to save the five people (and the death of the one person on the other track is just a tragic, foreseen side effect), switching the tracks may be permissible. 
However, if your intention is to actively harm the innocent person, it is not permissible.

\noindent\textbf{Victim-centered Deontology} focuses on the rights of the victim, and rules are designed to protect individuals from being used as a mere means to an end.
It is rights-based rather than duty-based.

\noindent For example, in the transplant case, It is wrong to harvest one person's organs to save five, because this directly violates that person's right not to be used, regardless of the overall outcome.
In the Fat Man trolley problem, because pushing the Fat Man off the bridge uses his body as a physical barrier to stop the trolley without his consent, it is strictly forbidden.

\subsection{Paradox of Deontology}
Victim-centered deontology claims that the foundation of morality is the rights of the victim (e.g., the right not to be killed, or the right not to be used as a mere means).
Can this theory be agent-neutral, i.e. everyone shares the exact same universal goal?

\noindent\textbf{Scenario}: A villain is about to murder five innocent people. You can stop him, but only if you murder one innocent person yourself.
The universal goal here would be "Minimize the total number of rights violations in the world".
An agent-neutral view says you must murder the one innocent person (1 violation is better than 5 violations).
The problem is that Deontology strictly forbids you from murdering an innocent person to achieve a good consequence. Therefore, Victim-Centered Deontology cannot be agent-neutral.

\noindent Victim-Centered Deontology must be Agent-Relative.
The personal goal would be "Ensure that I personally do not violate anyone's rights, regardless of what happens around me".
So we do not kill the one innocent person, and keep our hands clean at the expense of five innocent people dying.

\noindent The paradox is that a Victim-Centered theory was supposed to protect the victim's right not to be killed.
We protected the one victim from us, but allowed five other victims to die, just so that we don't break the rules.

\noindent Victim-Centered Deontology accidentally collapses into Agent-Centered Deontology.
It claims to care about the victim, but when push comes to shove, it actually prioritizes the moral purity of the agent over the lives of the victims.
The paradox happens because Victim-Centered theories enforce an Agent-Relative constraint: 'I must not violate a right.'
By focusing on keeping my hands clean, the theory forces me to passively watch five other victims die.

\subsection{Non-Absolute Deontology (Ellis’ Paper)}
Under traditional Deontology, moral rules are absolute. For example, you should never kill under any circumstance.
This leads to Moral Catastrophes.

\noindent Imagine this scenario: a hidden nuclear device is about to explode. Unless you (Person A) violate your duty and torture an innocent person (Person B), ten, a thousand, or a million innocent people will die in the blast.
Since traditional Deontology forces us to let the millions die, we look to a softer theory called Non-absolute Deontology or Threshold Deontology.

\noindent\textbf{Non-absolute Deontology} says that moral rules are absolute until the consequences of following them become completely disastrous.
Once the threshold is crossed, the absolute rule breaks, and you are allowed to look at the consequences (and do the "bad" thing).

Anthony Ellis attacks this idea of a threshold. He argued that no plausible version of deontology can successfully capture what morality requires of us.
\textbf{Hijacked Plane Example}: Imagine terrorists hijack a plane with 400 people on board, and the terrorists threaten to blow up the plane unless you murder someone.
Ellis breaks the choice down into a mathematical spectrum:
\begin{itemize}
    \item Act 1: Murdering 1 person to save 1.
    \item Act 2: Murdering 1 person to save 2.
    \item ...all the way to Act 400: Murdering 1 person to save 400.
\end{itemize}
A non-absolutist deontologist would say that it is wrong to murder unless the consequences of refraining are sufficiently bad. So, they would easily say that Act 1 is wrong, but Act 400 is right.
Ellis' counter is: At exactly what point does the murder become right?
Any threshold point chosen is completely arbitrary. (If you say 399 people, then why not 398 people?)

\subsection{Degree of Wrongness}
Non-absolute deontologists defend against Ellis by introducing the concept of Degree of Wrongness.
They argue that the threshold isn't arbitrary because it scales based on Degrees of Wrongness. Some actions are "more wrong" than others, which means they require a much higher threshold of catastrophe to justify breaking them.
\begin{itemize}
    \item Stealing a car has a relatively low degree of wrongness. You might be morally allowed to steal a car to drive someone to the hospital and save one life.
    \item Murder has a massive degree of wrongness. You might need to save a million lives to cross the threshold and justify a murder.
\end{itemize}
They try to prove that the threshold system has a logical, sliding scale rather than just a random cutoff point.

\vspace{1em}
\noindent\textbf{Legal Model Trap (Ellis's Critique)}
To justify "Degrees of Wrongness," some deontologists try to borrow the structure of the legal system (which categorizes crimes into degrees, like 1st-degree vs 2nd-degree murder).

\noindent Ellis's Objection: This fails because the legal system is a man-made, conventional construct. Legal cutoffs are inherently arbitrary decisions made by judges/politicians.

\noindent Deontology is based on Intrinsic Wrongness (objective, universal truths based on reason). You cannot use a man-made, arbitrary legal framework to solve a problem about absolute, intrinsic moral facts.
Doing so reduces universal morality to a mere human convention.

\subsection{Paradox of Relative Stringency}
Ellis attacks the idea of Degrees of Wrongness, by arguing that they triggered the Paradox of Relative Stringency.
The entire foundation of Deontology is that moral rules are absolute duties (Categorical Imperatives).
By definition, an absolute cannot have "degrees." You cannot have one absolute rule that is somehow "stricter" or "more absolute" than another absolute rule.
Further, if you put weights on rules and comparing those weights against the weights of the consequences, Deontology turns into Utilitarianism.

\noindent\textbf{Conclusion}: Non-absolute Deontology is a failed theory.

\end{document}
